%% Manuscript revision: 2026-10-12. Public-version status is recorded separately.
%%
%% Every measured number must trace to a committed result file in this
%% repository (results/...) or to a run directory listed in "Reproduction".
%% Conditional arithmetic is defined by the displayed vectors and conventions.
%%
%% Build: tectonic -X compile main.tex
%%
\documentclass[10pt,a4paper]{article}

\usepackage[T1]{fontenc}
\usepackage{lmodern}
\usepackage[margin=2.6cm]{geometry}
\usepackage{booktabs}
\usepackage{array}
\usepackage{url}
\usepackage{xurl}
\usepackage{microtype}
\usepackage{enumitem}
\usepackage{xcolor}
\usepackage{amsmath}
\setlength{\belowcaptionskip}{6pt}
\usepackage[numbers,sort&compress]{natbib}
\let\cite\citep
\usepackage[colorlinks=true,allcolors=blue]{hyperref}
\hypersetup{
  pdfauthor={Yufeiyang Chen},
  pdftitle={When the Verifier Cannot Decide: Tracing Math-Verify Failures Through an RLVR Reward Pipeline},
  pdfkeywords={reinforcement learning with verifiable rewards; reward verification; mathematical reasoning; software reliability; large language models}
}

\newcommand{\mv}{Math-Verify}

\title{When the Verifier Cannot Decide: Tracing {Math-Verify} Failures\\
Through an {RLVR} Reward Pipeline }
\author{Yufeiyang Chen\thanks{The author's primary contributions were
selecting the research topic and determining the research direction and
methodology. AI agents (Claude and Codex) provided substantial assistance
with implementation, experiment execution, code-based analysis and
manuscript drafting. The author is responsible for the manuscript and its
claims.}\\[2pt]
School of Cyber Science and Technology, Sun Yat-sen University, Shenzhen, China\\
Correspondence: \texttt{fei76161@gmail.com}}
\date{12 October 2026}

\begin{document}
\maketitle

\begin{abstract}
An unsuccessful mathematical verification can yield the same reward as
a wrong answer. Preserving the distinction requires more than catching
an exception at the scorer. We audit Math-Verify, its wrapper, a reward
manager and batch aggregation in a pinned verl configuration.
Persistent-worker injections reproduce a known sparse-metadata failure:
moving the same failed row changes the outcome from a lost flag to a
batch exception. We trace its interaction with the wrapper's flagging
contract and test row-preserving aggregation in 21 regression fixtures.
A pre-declared replay of 93,733 stored responses observes 54 library-event
rows; longer limits recover one zero-to-one verdict and leave 18 rows
undecided. The realised sample therefore permits one to nineteen
recoveries under the declared reference protocol. Separate replays
cover all 193,767 stored answers and a human-labelled answer set.
Matched shared-pool experiments connect natural heavy inputs to a batch
abort under diagnostic 60-second library and 30-second outer limits;
default-limit and serial controls complete. Together, the measurements
specify regression requirements at the consumer boundary and show how
timeout and recovery policies interact under concurrency. The replay
rates describe the fixed corpora and protocols, longer limits are not
ground truth, and no training effect is inferred.
\end{abstract}

\noindent\textbf{Keywords:} reinforcement learning with verifiable rewards;
reward verification; mathematical reasoning; software reliability; large language models.

%% ---------------------------------------------------------------------------
\section{Introduction}

RLVR replaces a learned reward model with a program that checks each answer
against a reference \cite{lambert2024tulu3,shao2024deepseekmath,deepseek2025r1}.
For mathematics the check is usually symbolic: the answer and the reference
are parsed into SymPy expressions and compared, and a library such as
\mv{} \cite{mathverify} wraps the parsing and comparison in time limits
because some expressions are expensive to evaluate. The reward is binary,
but the checking program has a third outcome that the binary reward cannot
express: it may fail to decide. A comparison can time out, a parser can
raise, the subprocess that runs the check can be killed, and a process
pool shared by many rows can break.

Training frameworks have repeatedly run into this third outcome. In verl \cite{sheng2024hybridflow}, contributors have reported that a
signal-based time limit that could not be armed in a Ray worker thread made
every MATH reward zero until the limit was disabled \cite{verl5635}, that
disabling it then let training hang \cite{verl5839}, and that a dead or stuck
scoring subprocess yields persistent zero rewards \cite{verl8011}. The pull request that addresses the last report
\cite{verl8015} rebuilds a broken pool, kills a stuck one, and propagates
persistent failures upward ``so the reward manager can flag the row''.
Separately, contributors have already reported sparse-metadata loss and
batch exceptions, and proposed union-key assembly in one or both batch
paths \cite{verl6545,verl7151}. These records establish known failure
mechanisms and candidate repairs. They do not by themselves give a rate
estimate for the fixed response corpus studied here, or establish how
the \#8015 wrapper's failure reporting behaves when composed with its
manager and aggregation paths.

We audit the reward pipeline as an end-to-end contract. A wrapper can
correctly report an error while aggregation removes that report; a larger
timeout can recover a verdict while also keeping costly work alive for
longer. Component-level tests therefore need to be connected to row-level
batch behaviour. We combine executable mechanism tests, fixed-response
replay and matched shared-pool controls to examine these relationships in
one pinned configuration. We ask:
\begin{enumerate}[label=\textbf{RQ\arabic*},leftmargin=*]
\item When the reward wrapper reports a failure, does that information
      survive the reward manager and batch aggregation?
      (\S\ref{sec:mechanism})
\item Which verifier events and limit-sensitive verdict changes occur in
      the stored responses, and how do they compare with representation
      effects in a separate labelled answer set?
      (\S\ref{sec:prevalence})
\item How do time limits and concurrent pool recovery interact when
      naturally heavy inputs enter a scoring batch?
      (\S\ref{sec:consequences})
\end{enumerate}
The pinned framework comparison (\S\ref{sec:frameworks}) and conditional
advantage calculations (\S\ref{sec:update}) put these measurements in
context; they are not additional full-framework experiments or learning
outcomes.

\paragraph{Contributions.}
\begin{enumerate}[nosep,leftmargin=*]
\item \textbf{A reproduced mechanism and an executable aggregation contract.}
  We trace an already reported sparse-schema failure through the real
  wrapper, manager and aggregation functions, using persistent-worker
  injections to test failure reporting. A bounded prototype based on the
  existing union-and-fill approach preserves the reward vector and aligned
  metadata in 21 fixtures, providing a concrete regression target
  (\S\ref{sec:mechanism}).
\item \textbf{A paired measurement of verifier decisions and unresolved cases.}
  A pre-declared replay covers 93,733 responses and carries unresolved
  rows into a conservative bound. Separate extensions cover all 193,767
  stored answers and 5,556 human-labelled extracted answers, distinguishing
  timeout-sensitive verdicts from representation-sensitive ones
  (\S\ref{sec:prevalence}).
\item \textbf{Controlled evidence for timeout--concurrency interactions.}
  Natural heavy inputs connect the recovery wrapper's previously reported
  reset interaction to a manager-level batch abort under relaxed diagnostic
  limits. Matched default-limit, serial and single-heavy controls isolate
  the tested conditions. A pinned framework audit and conditional
  advantage calculations explain the surrounding policies and their
  possible consequences (\S\ref{sec:consequences}--\ref{sec:update}).
\end{enumerate}

\paragraph{Scope.}
The experiments ran on one machine (Apple M2 Pro, 10 cores,
16\,GiB, macOS, Python 3.12.8) against pinned code; framework comparisons
are source audits and advantage examples are conditional calculations. Rates describe
\emph{this filtered dataset re-verified under the declared protocol}, not
natural rates during training. The reference judgement is the same verifier
with longer time limits, not ground truth. No model was trained.

\paragraph{Why three kinds of evidence.}
Injected failures test whether information survives a particular code
path; they do not estimate how often a verifier fails on model output.
Stored-response replays supply counts under a fixed scoring protocol,
including unresolved cases. Matched heavy-input scenes then vary timeout
and recovery policies to isolate their effects on a batch. Keeping these
roles separate prevents a demonstrated mechanism, a corpus frequency and
a diagnostic batch abort from being mistaken for the same result.

%% ---------------------------------------------------------------------------
\section{Background}
\label{sec:background}

\paragraph{Reward computation in verl.}
In the pinned verl revision, the default routing sends the MATH datasets to
verl's own \texttt{math\_reward} scorer, not to \mv{}; the \mv{} wrapper is
an optional scorer that a user selects with a custom reward function, as
verl's own tests do. Rewards
are computed by a reward manager. The default \texttt{naive} manager calls
the scoring function in a thread pool (via \texttt{run\_in\_executor}) and does
not catch exceptions; the \texttt{rate\_limited} manager limits concurrency
(default~1), and on a timeout or exception records a reward of $0.0$ together
with an extra \texttt{timeout} or \texttt{error} key. The default reward
configuration starts 8 reward workers. When rewards are computed during
generation, each response is sent to a randomly chosen worker; when they are
computed after generation, the batch is split into contiguous chunks.
Routing and worker/chunk counts were checked at the tested head
\texttt{0c4c3d4}; they do not imply eight request threads per worker.

\paragraph{The \mv{} wrapper in pull request \#8015.}
The wrapper at head \texttt{0c4c3d4} keeps a module-level
\texttt{ProcessPoolExecutor} with at most four spawned workers. It submits
each row and waits at most 30\,s. On an outer timeout it cancels its own
task, kills every worker of the current pool, and returns
\texttt{timeout\_score} (default $0$). On the first broken-pool error it
rebuilds the pool and retries once; a second consecutive failure is raised
to the caller. Python's executor marks every pending future of a broken pool
as failed, so one dying worker affects every row in flight on that pool.

\paragraph{Time limits inside \mv{}.}
\mv{} 0.9.0 enforces its own limits with \texttt{signal.alarm}: 5\,s for
parsing each side and 5\,s for each candidate pair in the comparison. With
\texttt{raise\_on\_error} left at its default (false), a parse that times out or raises
returns an empty list and a pair that times out or raises counts as
``not equal''; if another pair (for instance the plain-string pair) matches,
the overall verdict is still true. The library's timeout exception derives
from \texttt{BaseException}. Because \texttt{signal.alarm} only works on
the main thread, calling it from another thread raises an error, which the
library rewrites into an explicit message; the wrapper calls the library in
the main thread of its pool subprocesses.

%% ---------------------------------------------------------------------------
\section{Verdicts, failures and contracts}
\label{sec:contracts}

We distinguish four kinds of zero (Table~\ref{tab:events}). This study traces
undecided verdicts and reference-side failures through the scoring pipeline.

\begin{table}[h]
\centering
\caption{What a zero reward can mean.}
\label{tab:events}
\small
\begin{tabular}{@{}p{3.4cm}p{6.3cm}l@{}}
\toprule
Event & Example & Responsibility \\
\midrule
Ordinary negative & declared unequal, or no answer admitted by the extraction policy & policy \\
Policy over budget & truncated response, exhausted turns & policy \\
Verifier undecided & library timeout or exception, dead worker, broken pool & infrastructure \\
Data-side failure & the reference answer itself cannot be parsed & data \\
\bottomrule
\end{tabular}
\end{table}

Whether a zero for an undecided verdict is a \emph{defect} depends on the
contract the pipeline claims to implement. Under a \emph{budget contract}
(``correct means provably correct within the compute budget'') a timeout
scored as wrong is intended. An \emph{observability contract} instead
requires a failed attempt to remain distinguishable from a negative verdict,
even when both receive zero. A separate \emph{exclusion policy} must specify
whether an undecided row is removed from the group statistics, the loss,
and its denominator. A promise to flag a failure does not by itself promise
to exclude it from optimisation. We call
behaviour a defect only when a pinned implementation states a contract in
code comments, documentation or a pull-request description and its runtime
behaviour departs from it. Where no contract is stated we report a
\emph{contract ambiguity}. We also separate the \emph{default protocol}
(library defaults, framework routing and platform as shipped) from
\emph{diagnostic protocols} that change parameters to make failures visible;
results under the latter are never reported as default behaviour.
The uniform-key expectation stated in the discussion of \#6545
\cite{verl6545} is relevant here: our question is how that assumption
composes with a manager that emits failure-only keys. We test the
specified returned reward and metadata views, rather than assuming that
every sparse dictionary is valid under every consumer's contract.

%% ---------------------------------------------------------------------------
\section{What happens to a failure flag}
\label{sec:mechanism}

Pull request \#8015 states an observability contract for persistent
failures: the wrapper raises ``so the reward manager can flag the row'', and
the \texttt{rate\_limited} manager does write an \texttt{error} key. We tested
whether the flag reaches the consumer.

\paragraph{Setup.}
We ran verl at the pull-request head with the real \texttt{rate\_limited}
manager, the real reward-loop workers and both aggregation functions (the
per-sample path used during generation and the batch path used after it).
Generation was replaced by fixed response texts, so no model was loaded.
Persistent failure was injected by making the scoring subprocess call
\texttt{os.\_exit(17)} on a marked response, on every attempt; a ledger
recorded each injection. Thirteen cases varied the position of the failing
row, whether whole chunks or whole batches failed, and the manager.
The two recorded runs are in \texttt{results/c1c/run1/} and
\texttt{results/c1c/run2/}, against execution commit \texttt{cc87848}.
They expose the same outcomes on repeated execution of the declared cases.

\paragraph{Results.}
Both aggregation paths take the set of extra-information keys from the
\emph{first} row of a chunk and then index every row by those keys.
This is the sparse-schema mechanism already described in
\cite{verl6545,verl7151}. Our failure injection makes the incompatible
manager outputs and their returned failure information explicit:
\begin{itemize}[leftmargin=*]
\item If the failing row is not first, the batch completes, the failing row
      contributes $0$, and the reward returned by \texttt{extract\_reward}
      carries no failure flag: from that return value a failed scoring is
      indistinguishable from a wrong answer. (The raw per-sample column on
      the generation path still holds the flag; we did not run a full trainer
      to see whether anything reads it.)
\item If the failing row is first and other rows are normal, aggregation
      raises \texttt{KeyError} and the batch aborts. Moving the \emph{same}
      originally correct answer from a later position to the first reproduced
      the switch, so position, not content, decides between the two outcomes.
\item If an entire chunk fails, mismatched keys or metadata across chunks
      trigger assertion failures; if the entire batch fails, zeros and flags
      are returned consistently.
\item Under the default \texttt{naive} manager the same failure aborts the
      batch with the true cause, because the manager does not catch it. The
      flagging design therefore only applies with \texttt{rate\_limited}.
\end{itemize}

\begin{figure}[htbp]
\centering\small
\fbox{\begin{minipage}{0.94\textwidth}
\centering
\textbf{Wrapper exception $\longrightarrow$ manager returns reward 0 and an error flag}

\medskip
\begin{tabular}{@{}p{2.5cm}ccc p{5.6cm}@{}}
\toprule
Failed row & Row 1 & Row 2 & Row 3 & First-row-key aggregation \\
\midrule
Row 2: rewards & 1 & 0 & 1 & first row has no error key \\
Error metadata & --- & error & --- & error column is omitted \\
\addlinespace
Row 1: rewards & 0 & 1 & 1 & first row contains the error key \\
Error metadata & error & --- & --- & indexing a normal row raises \texttt{KeyError} \\
\bottomrule
\end{tabular}

\medskip
\textbf{Union-and-fill aggregation: retain the key and align missing entries explicitly}
\end{minipage}}
\caption{Schematic of the row-position effect. Common metadata is omitted;
three rows illustrate the first-row key rule rather than a new measured
batch. Both real aggregation paths exhibit this effect. The tested repair
prototype applies only to batch assembly after scoring, with missing
entries represented by \texttt{None}.}
\label{fig:failure-flag}
\end{figure}

\paragraph{Row-preserving aggregation.}
Union-key assembly with explicit missing entries was proposed upstream
\cite{verl6545,verl7151}; it is not a new repair algorithm. We use a
bounded local version to specify and test the failure-recording contract.
Let a batch contain ordered pairs $(r_i,e_i)$, where $r_i$ is the numerical
reward and $e_i$ is the row's optional metadata dictionary. An aggregation
that preserves failure information must satisfy three requirements:
\begin{enumerate}[nosep,leftmargin=*]
\item \emph{Key completeness:} retain every key in
      $K=\bigcup_i \mathrm{keys}(e_i)$, including keys absent from the first row.
\item \emph{Row alignment:} each metadata column has one entry per returned
      reward, in the same order after any padding and trimming. Absence
      requires an explicit representation rather than dropping the column.
\item \emph{Reward preservation:} adding observational metadata leaves
      the numerical reward vector unchanged; exclusion from an update is
      a separate policy decision.
\end{enumerate}
The first-row key rule violates completeness, and direct indexing assumes
every row has every retained key. The fixtures below exercise these two
failure modes at each row position and check alignment across worker
chunks. These are aggregation requirements; they cannot expose a library
failure that was never included in $e_i$.

\paragraph{Row-permutation equivariance.}
For a batch without padding, the same requirements imply an executable
metamorphic check: reordering fixed scorer outputs should only reorder
the returned rewards and metadata. If $\pi$ is a row permutation and
$A(B)$ assembles batch $B$, then for every retained metadata key $k$,
\[
 \mathrm{scores}(A(\pi B))=\pi\,\mathrm{scores}(A(B)),\qquad
 \mathrm{column}_k(A(\pi B))=\pi\,\mathrm{column}_k(A(B)).
\]
The retained key set should be unchanged. Moving a flagged row to the
first position violates this requirement in the original code even
though its scorer output has not changed. The eight error-position and
eight timeout-position fixtures below exercise the same row-alignment
requirement for the local prototype.

\paragraph{A bounded aggregation repair.}
We tested a local prototype on the pinned batch aggregation method, changing
only two assignments: take the union of extra keys over all rows, and fill
missing entries with \texttt{None} in object-valued columns. Twenty-one
fixed scorer-output fixtures cover an \texttt{error} or \texttt{timeout}
at each of eight positions, mixed flags, empty extras, normal outputs and
padding/trimming across two workers. The original method omits columns in
16 fixtures and raises in three; the prototype returns the expected reward
vector and aligned extra values in all 21. Upstream source was unchanged.
This is a local regression of batch assembly with fixed scoring outputs and
substituted Ray transport, not a verifier or training run. It does not repair
the generation-path concatenation, naive-manager exceptions or silent
library failures, and downstream acceptance of \texttt{None} or a masking
policy remains untested.

\paragraph{Library-internal timeouts do not reach the wrapper's timeout branch.}
The wrapper maps the library's timeout exception to \texttt{timeout\_score}.
With default parameters the library catches its own timeouts first, so the
branch is not reached: setting \texttt{timeout\_score} to $-7$ left
library-level failures at $0$ or $1$, and only the outer 30-second wait and an
explicitly raised timeout returned $-7$. In a controlled slow comparison, two
originally correct answers dropped to $0$, while a third still scored $1$
through the string fallback. Because no contract states which timeout layers
\texttt{timeout\_score} is meant to cover, we report this as a contract
ambiguity, not a defect.

\paragraph{Mechanism evidence and replay incidence.}
The aggregation behaviour is a property of real verl code at the tested head,
exercised end to end through the manager. The failure itself was injected; it
does not show that natural inputs reach the propagation path. The next
section measures what natural inputs do.

%% ---------------------------------------------------------------------------
\section{How often: re-verifying 93,733 responses}
\label{sec:prevalence}

\subsection{Protocol}

The sample, replay settings, observation checks, follow-up selection and
decision rule were fixed before the main sample was scored (execution
snapshot \texttt{9a88991}). The sample map and threshold calculations are
distributed under \texttt{plan/} in the companion archive. The finite-sample
recovery range below is an additional interpretation of the same frozen
counts, not a change to the pre-declared rule.

\paragraph{Sample.}
OpenR1-Math-220k \cite{openr1math} (configuration \texttt{default}) contains
93,733 problems from NuminaMath~1.5, each with one to six reasoning traces
generated by DeepSeek-R1; according to the dataset card, traces were verified
with \mv{} or, for about 12\% of samples, an LLM judge, and every problem has
at least one trace judged correct. For each
problem we chose one response uniformly with a fixed seed, using a hash of
the problem text (the dataset's \texttt{uuid} is missing for 664 rows and
shared by 60 pairs of different problems). The full response text, including
the reasoning, was passed to the wrapper with the dataset's \texttt{answer}
as reference. The dataset was already filtered by the same library, so
nothing about the direction of bias can be assumed.

\paragraph{Isolated replay.}
Each measurement process owns one wrapper and its pool, submits one row,
waits for it, and only then submits the next; a new pool is first warmed with
a trivial row whose result is discarded. There is therefore no queueing and
no cross-row interference by design. The wrapper code, the library limits
(5\,s), the outer limit (30\,s) and \texttt{timeout\_score}$=0$ are unchanged.

\paragraph{Observation layer.}
In the parent we wrap pool creation, task submission and each wait, recording
outcomes, pool generations, resets and every kill. In the child, before the
unchanged worker function runs, we wrap \mv{}'s \texttt{parse} and
\texttt{verify} and record each step (reference parse, answer parse, each
comparison pair) with wall and CPU time, the result type, library log
records, and the arrival of every alarm signal. Records carry sequence
numbers; a synthesis step cross-checks parent and child records row by row
and marks any row whose records are incomplete or inconsistent as undecided.
Turning the layer off and on changed no score on 500 pilot rows. Eighteen
negative controls, each corrupting one record type, were all detected.
The self-test records are in \texttt{results/s2\_rev2/selftest/report.json}.

\paragraph{Reference judgement and follow-ups.}
Every row scored $0$ by default was re-scored with the library limits raised
to 60\,s and the outer limit to 300\,s (the \emph{reference}). A pre-drawn
random set of 1,000 default-$1$ rows was also re-scored to check for
one-to-zero changes in those sampled rows. Rows with a default or reference event,
a score change, or a limit-sensitive change were replayed three times under
the default setting. Rows with a library timeout or any step over 1\,s, plus a
pre-drawn random 1\%, were re-scored with library limits of 2\,s and 10\,s.
Another pre-drawn 1\% of rows were re-run unchanged to measure test--retest
instability.

\paragraph{Metric and decision rule.}
Let $k$ be the number of rows scored $0$ by default and $1$ by the reference,
and $u$ the number left undecided (the reference itself timed out, the row was
interrupted twice, or the local resource guard stopped it). The denominator is
always $n=93{,}733$. For the realised sample, the possible zero-to-one
recovery share under this reference protocol is partially identified:
\[
  \frac{k}{n}\ \leq\ q_{\mathrm{recover}}\ \leq\ \frac{k+u}{n}.
\]
These finite-sample bounds assign unresolved rows to either endpoint and
need no assumption of independent, identically distributed problems.
They concern score recovery under the specified verifier, not mathematical
correctness. The pre-declared decision rule additionally uses binomial-model
Clopper--Pearson bounds \cite{clopper1934}. These are working-model bounds
for the declared rule, not distribution-free bounds for heterogeneous
problems or guarantees about a training population.
Before the run we fixed: continue to the consequence stage if the one-sided
95\% lower bound on $k/n$ is at least $10^{-3}$ (here $k\ge111$), or if a
natural input reproducibly raises through the wrapper, or if an outer timeout
spills over to other rows; stop this line if none holds and the one-sided
95\% upper bound on $(k+u)/n$ is below $10^{-4}$ (here $k+u\le4$); otherwise
report the result as \emph{in between}.

\subsection{Results}

\begin{table}[h]
\centering
\caption{Default-protocol outcomes on the main sample ($n=93{,}733$).}
\label{tab:default}
\small
\begin{tabular}{@{}lr@{}}
\toprule
Outcome & Rows \\
\midrule
Score 1 & 61,246 \\
Score 0, no verifier event & 32,442 \\
Score 0, comparison timed out (library) & 27 \\
Score 0, comparison raised (library) & 12 \\
Score 0, answer parse raised (library) & 4 \\
Score 0, answer parse timed out (library) & 1 \\
Score 0, reference parse timed out (library) & 1 \\
Outer timeout, broken pool, exception raised by the wrapper & 0 \\
\bottomrule
\end{tabular}
\end{table}

\paragraph{Events are rare; selected replays are stable.}
Under the default protocol 45 zero-scored rows had a verifier-side event
(Table~\ref{tab:default}); a further 9 rows had a comparison error in one
candidate pair but still scored $1$ through another pair. No row hit the outer 30-second limit, broke the pool, or made the wrapper
raise, and no row took more than 10\,s. All 47 rows selected for replay gave
the same outcome three times out of three, and the 963-row retest changed
nothing.

\paragraph{One score change, eighteen undecided.}
The reference turned exactly one default $0$ into a $1$ ($k=1$). Eighteen rows
stayed undecided: fourteen in which the 60-second comparison also timed out,
one in which parsing the reference answer timed out, two interrupted twice and
one stopped by the local resource guard. The last three were interrupted or
stopped by the resource guard then in use on this 16\,GiB machine (since
retired); they are local interventions, not natural outcomes. All 1,000 default-$1$ rows stayed $1$.
The finite-sample recovery interval is
$[1.07\times10^{-5},\,2.03\times10^{-4}]$. With $k=1$, $u=18$ the
one-sided 95\% binomial-model upper bound on $(k+u)/n$ is
$2.97\times10^{-4}$; the two-sided 95\% interval for $k/n$ is
$[2.7\times10^{-7}, 5.9\times10^{-5}]$. Neither threshold is met, so the
pre-declared verdict is \emph{in between}. Of 928 rows re-scored with 2\,s
and 10\,s limits, three changed with the limit: two scored $0$ at 2\,s and $1$
at 5\,s and 10\,s, and one (the $k$ row) scored $1$ only at 10\,s.

The bound rests on the reference not being backfilled. Had all eighteen
undecided rows scored $0$ under some longer limit, the upper bound would fall
below $10^{-4}$; had all scored $1$, $k$ would be 19, still far from the
continue threshold. The supplementary diagnostics below were specified in
advance for exactly these rows and are reported separately; they do not
change the formal result.

\subsection{The one change, measured again}

The single $k$ row changes from $0$ to $1$ when the comparison limit is
relaxed; this score change is not an independent proof of mathematical
correctness. Re-run ten times in sequence on one
warmed pool, it timed out at the 5\,s alarm every time (comparison wall
5.002--5.005\,s) and scored $0$; under the reference it scored $1$ every time
with comparison wall 5.82--7.51\,s (median 6.21\,s). The ten repeats share a
pool and SymPy's caches and are not independent samples, and we have no
measurement on other hardware; all ten repeats show the same limit-dependent
pattern on this machine. The first repeat had a gold-parse time of 0.11\,s
against $<$0.001\,s afterwards; cache state was not controlled. It behaved
the same way: $0$ at the 5.005\,s alarm
under the default, and $1$ under the reference with the slowest comparison of
the ten (7.51\,s).

\subsection{Heavy natural inputs}
\label{sec:prevalence-heavy}

Four undecided rows were singled out because the reference either could not
finish them or stopped them for memory. We ran each in a fresh process that
calls the wrapper's worker function directly with the 60-second library
limits, dumped all thread stacks every 10\,s, and stopped a run at 1.5\,GiB of
memory footprint or 120\,s of wall time after a final stack dump. These are
stack-capture endpoints, not resource limits of the method.

All four stacks end in the same place: \mv{}'s numeric equality calls SymPy's
\texttt{doit()}, which evaluates an integer power in
\texttt{Integer.\_eval\_power}. Three rows express the same recursive-function
task for $f(4,1981)$, with identical recurrence rules. Their \emph{reference answers} contain nested
powers: $-3+2^{2^{1984}}$, $2^{2^{1984}}-3$ and $2^{2^{1984}-3}$.
The last is not the same expression as the first two; we report the corpus
strings rather than settling their mathematical correctness. The responses
all answer $2\uparrow\uparrow 1984-3$. In each, footprint grew at 82--88\,MiB/s over the
first 15\,s and passed the 1.5\,GiB endpoint within about 18\,s of process
start, with no sign of levelling off within that window; a truncated
trace says nothing about whether or when the computation would end. The fourth row's reference answer is $0$ and its response
contains $5^{5^{5^{5^{5}}}}$; there memory grew slowly (peak 351\,MiB) and the
60-second alarm stopped the comparison on time, scoring $0$, as in the main
measurement. Under the default 5-second limit all four were stopped by the
alarm and scored $0$; none reached the outer limit or harmed the pool.

So in the default protocol the library's alarm is what contains these
inputs: they score $0$ in about 5\,s and leave the pool intact. Whether the
0 is wrong is not settled here; all four rows remain undecided. The trigger can sit in the reference answer, in which case any
response to that problem whose comparison takes the same numeric path would
meet it --- in GRPO, possibly several members of one group. The corpus
extension below counts within-question events under isolated replay; it does
not reproduce simultaneous generation or training. The shared-pool experiment
below tests selected inputs at explicitly relaxed limits; disabled limits
are only a source-level observation.

\subsection{The remaining fourteen}
The other fourteen undecided rows were re-scored once each with 300\,s
library limits (outer limit 1,500\,s), a bounded sensitivity analysis rather
than a ground truth. None scored $1$. Eleven comparisons reached their
300\,s wall-clock limit and one reference-answer parse reached that limit;
the limit applies to each library operation, not to the entire response.
Two comparisons
finished and returned false, after 81\,s and 281\,s. Counting the four heavy
rows as still undecided, this gives $k_{300}=1$ and $u_{300}=16$: the one-sided
95\% upper bound on $(k_{300}+u_{300})/n$ is $2.72\times10^{-4}$ and the
lower bound on $k_{300}/n$ is $5.5\times10^{-7}$, so the analysis stays between
the two pre-declared thresholds, as the formal result does. A preliminary
reading of the answer strings, which is not ground truth, judged ten of the twelve rows that
remain undecided to state the reference answer and the other two not to; it
agrees with both finished zeros (one reading ``different'', one with nothing
to judge). A limit sixty times the default thus decided two of the fourteen,
both as zeros.

The run was interrupted once after seven rows and resumed
from the same directory; the eighth row was re-run from the start. During the
first part the machine carried heavy unrelated load (median one-minute load
average 10.5), and one memory-pressure reading was critical, which paused
dispatch for 32\,s without stopping the row in progress (the one that
finished at 281\,s, with 208\,s of comparison CPU time). The eleven
timed-out comparisons used approximately 250--299\,s of CPU time. These
results show that the stated wall-clock protocol did not finish those
operations on this host; they do not establish a hardware-independent
minimum computation time.

\subsection{Extension to all stored long-form answers}
\label{sec:allanswers}

After the main result was fixed, we specified a replay of only the 100,034
answers outside the one-answer-per-question sample. Question identity is
the exact text's SHA-256, with an answer index appended; the 93,733 question
strings contain one to six answers each, 193,767 in total. The original
sample, its formal report and decision thresholds were not re-scored or
changed. Each new answer used the same isolated 5\,s / 30\,s protocol;
all 32,704 answers not scoring one were re-run at 60\,s / 300\,s, along with
1,000 default-one answers selected by a fixed hash ordering of question
key and answer index. Code, data, selection and baseline hashes were locked
before replay.

The complement contains one default-zero/reference-one change and fourteen
undecided answers (thirteen comparison timeouts and one reference-parse
timeout). All 1,000 controls remain one. Joining these \emph{counts} to the
audited main counts gives two changes and 32 undecided answers in 193,767:
the corresponding finite-corpus proportions are $2/n$ and $34/n$, or
0.00103\% and 0.01755\%. These are descriptive endpoints, not confidence
bounds, and do not apply the main sample's decision rule. The second change
is another answer to the main sample's one limit-sensitive question;
it is not a second distinct problem exhibiting that effect.

The full default replay has 103 answers with an observed library failure
event; 85 score zero and eighteen still score one through another accepted
comparison. Of the 92,287 questions with at least two answers, 59 have an
observed event and 42 have events on two answers. Define an ordered pair
by choosing an observed-event answer and one other stored answer to the
same exact question. There are 103 such pairs, of which 84 have events at
both endpoints (81.55\%); no other endpoint has incomplete observation.
This denominator excludes single-answer questions and weights questions
by the number of event anchors and other answers. All 42 multi-event
questions have two stored answers. The same reference and question content,
different replay times and uncontrolled caches can jointly explain
co-occurrence. We do not estimate a correlation parameter, a training
failure rate or dependence in eight-answer GRPO groups from these counts.

The default complement replay finished in 14.2 minutes. Its reference
controller disappeared once with one unfinished answer; the reason was
not recorded. The same directory and version lock were resumed, retaining
completed work and the interrupted attempt. The reported 20.4-minute
reference duration describes the resumed controller only, not all reference
work or the intervening gap. No red-pressure intervention occurred in this
extension (maximum recorded pressure was warning). Raw-event reassembly
and a separate count/pair reconstruction match the stored report.

\subsection{Cross-check with human-labelled extracted answers}
\label{sec:labels}

We used the released human-labelled 1.5B-model answers of
Xin~\cite{xin2026groupcorr,humanlabels}: 5,556 answers, with 2,822 positive
and 2,734 negative labels, drawn from 3,041 groups. The author's sampling
code stratifies by category and verifier opinions; this is not a random
sample of a training population. Only 72 groups contain all eight answer
positions. Two annotation columns agree on the 200 double-labelled rows;
we did not independently adjudicate the labels. The input is an extracted
answer, not the full generated reasoning used in the main sample.

Each answer was replayed in isolation through the same pinned wrapper,
first verbatim and then, if it did not already contain \verb|\boxed|,
enclosed in \verb|\boxed{...}|. Each presentation was scored with default
5\,s library / 30\,s outer limits and with 60\,s / 300\,s limits. The 24
already-boxed answers reuse their recorded scores rather than creating a
second identical trial. Thus there are 22,176 actual trials and 5,556
label-aligned rows per presentation and limit condition. Input strings,
source hashes and alias mappings were fixed before measurement. All rows
have complete observation records, and no resource intervention occurred.
An ordinary zero here means no observed failure event, not necessarily a
completed unequal comparison: an empty extraction under the specified
format policy also produces zero.

\begin{table}[h]
\centering\small
\begin{tabular}{@{}lrrrr@{}}
\toprule
Presentation & Label-positive / score 0 & Label-negative / score 1 & Disagreement & Event zeros \\
\midrule
Verbatim & 1,487 & 43 & 1,530 & 3 / 2 \\
Explicit delimiter & 113 & 91 & 204 & 4 / 3 \\
\bottomrule
\end{tabular}
\caption{Label disagreement at default and extended limits: the scores are
identical between the two limits. The last column gives default / extended
failure-event counts; all those answers have negative human labels.}
\label{tab:labels}
\end{table}

Under either presentation, no observed failure-event zero intersects a
positive human label (0/2,822 label-positive answers). Extending the limit
eliminates one comparison-timeout event per presentation but changes no
score. This is a zero result on the selected sample; it is not a proof that
timeouts cannot hide correct answers elsewhere. The presentation change
turns 1,431 zeros into ones and nine ones into zeros, reducing disagreement
but increasing label-negative ones from 43 to 91. Most disagreement
therefore cannot be attributed to the observed timeout events. The delimiter
condition is an explicit input transformation, not a repair demonstrated
on full generated responses or a new verifier method. The underlying label
study already analyses verifier disagreement and group dependence; our
addition here is the failure trace under the specified execution protocols.

%% ---------------------------------------------------------------------------
\section{Consequences of heavy natural inputs}
\label{sec:consequences}
\subsection{Matched single-heavy batches}

The two unchanged wrappers were compared at \texttt{0c4c3d4} and
\texttt{acad987}, with a process-pool initializer setting the library limits
inside each spawned child's main thread. The original scoring functions,
manager and aggregation functions were executed; fixed text, an index
text-tokenizer and a local substitute for Ray transport replaced generation
and remote dispatch. One naive manager used eight request threads and
64 rows. This is a controlled amplification at one reward actor, not the
full default split across eight actors. Four extra score-one answers warmed
each fresh scene sequentially; this does not guarantee that all four lazy
pool children had been warmed.
The manager is instantiated afresh for each batch, while the scene retains
its event loop, thread executor and wrapper pool. Following batches therefore
probe scorer/pool state, not survival of the same manager object.

The matrix comprises two wrappers, two library limits (unmodified 5\,s or
60\,s), and fixed 30\,s outer waits. The 64 controls were selected by a
fixed hash ordering from the previous score-one set, independently of these
new outcomes. For each wrapper/limit pair, an all-control batch ran once;
the response-heavy input at position 1 or 33 and one reference-heavy
input at position 1 each ran three times. Every scene ended with a following
64-control batch in the same pool. The 40 scenes had no local resource
intervention; maximum recorded pressure was normal. Synthetic tests first
checked on/off score equality, default versus relaxed library timeouts and
the two outer-timeout pool policies. The natural scenes used the specified
30\,s waits, not the synthetic test's shortened wait.

Every heavy row scored zero, all 63 unaffected controls scored one, and all
following controls scored one. Under the 5\,s library limit, the heavy
scoring calls returned after 5.03--5.34\,s with no pool reset. Under the
60\,s library / 30\,s outer diagnostic, the heavy call returned zero at
30.00--30.18\,s in both wrappers. The pull-request wrapper reset the pool
in all nine heavy scenes at this diagnostic setting; the main wrapper reset
it in none. In main, the child's heavy computation ended approximately
29.72--29.92\,s \emph{after} the following control batch had finished,
showing that returning zero did not stop the background work. Its scene
shutdown waited for that work; it was not an observed training stall.

No shared-pool score damage was observed with \emph{one} heavy row in
these particular batches. This is not evidence that arbitrary concurrent
inputs are safe. More controls could consume the remaining pool capacity,
and the default concurrent protocol could reach the outer deadline by
queueing or cumulative comparison work. The diagnostic distinguishes
cleanup policies, not the deployment frequency of a default fault.
\subsection{Three-heavy, eight-row manager controls}

A separately specified matrix used the first eight of the same control
answers. For each wrapper and library limit, both a naive manager with
eight request threads and a rate-limited manager at its default concurrency
one scored either a pure-control batch or a batch with the three
reference-heavy rows first. Every scene then scored eight controls. All
16 scenes completed without resource intervention. This remains one
locally invoked reward actor; it does not execute a full Ray trainer.

\begin{table}[h]
\centering\small
\begin{tabular}{@{}llp{0.27\linewidth}p{0.27\linewidth}@{}}
\toprule
Library limit & Manager & Recovery wrapper & Main wrapper \\
\midrule
5\,s & Naive & Three zeros, five ones & Three zeros, five ones \\
5\,s & Rate-limited (1) & Three zeros, five ones & Three zeros, five ones \\
60\,s & Naive & BrokenProcessPool; no batch return & Three zeros, five ones \\
60\,s & Rate-limited (1) & Three zeros, five ones & Three zeros, five ones \\
\bottomrule
\end{tabular}
\caption{Three natural reference-heavy inputs in an eight-row batch,
with 30\,s outer waits throughout. Pure controls and following controls
all score one. The 60\,s condition is a diagnostic, not a library default.}
\label{tab:threeheavy}
\end{table}

The pull request's public discussion had already reported a second
\texttt{BrokenProcessPool} under concurrent resets with a minimal verifier
stand-in, and its author explicitly retained that propagation
\cite{verl8015discussion}. Our experiment extends that engineering record
with real verifier inputs and the manager's batch outcome; it does not claim
to discover the reset interaction.
The diagnostic naive recovery-wrapper scene raised
\texttt{BrokenProcessPool} after repeated pool resets; its five control
calls individually returned one, but no aggregate batch reward was
returned. The traceback reaches the wrapper's persistent broken-pool
raise. Two additional traced repetitions reproduced the batch abort
(three traced runs in total). A fourth run disabled child step/DEBUG
observation and also aborted; parent scheduling logs and the time-limit
initializer remained in place, so this is not a completely uninstrumented
run. The four matched main-wrapper scenes returned three zeros and five
ones. No subprocess-exit fault was injected, and no red-pressure guard
intervened. Default-limit and serial controls in the initial matrix each
ran once; we do not assign them a frequency estimate.

On a batch exception, the harness records it, drains requests already
started by the real gather, and only then issues the controlled following
batch. Those batches all return eight ones. This demonstrates recovery
under that harness, not that an actual training loop automatically resumes.
The contrast is a consequence of timeout/recovery policy interacting with
concurrency: it is not evidence that the declared zero-on-timeout policy is
itself a defect, nor that the diagnostic abort occurs at default limits.
It also does not establish deployment incidence or satisfy the original
main-sample decision rule retrospectively. The formal result and thresholds
remain unchanged.

\subsection{Sequential parent calls and a longer outer wait}

Fourteen additional, separately specified scenes reused the first eight
controls, the three reference-heavy inputs and the four warm-up answers.
Each wrapper had seven one-time scenes: pure-control and three-heavy inputs
under direct sequential calls at 5\,s / 30\,s and 60\,s / 30\,s; and
pure-control, single-heavy and three-heavy inputs under the eight-thread
naive manager at 60\,s / 300\,s. Each scene ended with eight controls.
All fourteen completed, with no resource intervention and maximum pressure
normal. Every normal answer scored one, and every heavy answer scored zero.

Under sequential parent calls at 60\,s / 30\,s, the recovery wrapper reset
its pool three times but returned all row results. Main made no reset;
its child's work continued after each outer return. The recorded child
intervals overlap even though the parent invokes the wrapper sequentially,
so this control does not imply entirely serial background computation.
At 60\,s / 300\,s the library's comparison deadlines ended the heavy work
before the outer waits; neither wrapper reset its pool, both naive batches
returned, and all following controls scored one. These controls are
consistent with the earlier batch abort arising from concurrent
outer-timeout recovery, rather than an isolated heavy call. They do not
establish that every concurrency pattern or timing environment behaves
the same way.

\section{Other frameworks}
\label{sec:frameworks}
Scoring paths were inspected at fixed revisions. The function-body and
string checks below test small parts of those paths; full deployment
configurations were not executed.

\paragraph{TRL.} TRL's built-in \texttt{accuracy\_reward}~\citep{trl} calls
\mv{} in the calling thread, with no worker process. On the main thread it
passes a 10\,s parsing and a 5\,s comparison limit; off the main thread it
passes \texttt{None} for both, which disables the alarm, and lowers \mv{}'s
log level to hide the resulting warnings. An unparseable reference answer
yields \texttt{None}, which the GRPO trainer turns into NaN and so keeps as a
flag; a comparison that times out returns false and becomes an unflagged $0$.
The GRPO trainer calls synchronous reward functions directly in its training
loop, without exception handling. Under an ordinary script launch that loop
is the main thread, so the limits apply.

\paragraph{Ray.} The Ray Train example for TRL GRPO~\citep{raytrl} copies
\texttt{accuracy\_reward} and sets both limits to $0$, stating that the
default limits are incompatible with Ray. In \mv{}~0.9.0 a limit of $0$ or
\texttt{None} means no alarm at all. A run that follows this example therefore
has no library deadline if a comparison reaches the integer powers of
\S\ref{sec:prevalence-heavy}; the verifier is called inside the training
process. This is a statement about the code; we have not run such a
configuration.

\paragraph{verl main.} At the main-branch revision we read (\texttt{acad987}),
verl's default scorer~\cite{verlmain} routes by data source and does not use \mv{}: the
MATH sources go to a rule-based checker, with \mv{} offered only in a comment
as a manual override; DAPO and AIME sources go to another rule-based checker;
six NuminaMath sources go to \texttt{prime\_math}. Its first symbolic path
has an intended 10\,s subprocess timeout and catches exceptions as false;
we have not executed that path across multiprocessing start methods.
The fallback calls \texttt{math\_equal(..., timeout=True)}, passing a
boolean to a duration parameter. More directly, its \texttt{symbolic\_equal}
uses \texttt{with timeout\_limit(...)} although that helper returns a
decorator function. Each such block raises \texttt{TypeError} before parsing
or comparison, and the exceptions are swallowed. In a bounded check using
the fixed function bodies and call sentinels, the fallback returned false
without reaching any symbolic operation. Earlier string or numeric branches
can still return true; this is not a claim that all prime\_math scores are
zero. The \#8015 wrapper we test therefore sits on the
opt-in \mv{} route, and our results apply to configurations that choose it.

\paragraph{OpenRLHF.} OpenRLHF~\citep{openrlhf} does not use \mv{}. Its math
grader (adapted, per its header, from the OpenRLHF-Agent project) compares
with SymPy \texttt{simplify}
but first refuses to hand SymPy any expression containing \verb|^{| or
\verb|^(|, a chained power, an exponent of two or more digits, or more than two
unknown letters; such pairs are graded unequal without evaluation. On the main
thread \texttt{simplify} runs under a 120\,s alarm; off the main thread (the
code comment names Ray actors) it runs without a limit. A timeout or other
exception is graded false, again without a flag. This is a contrasting input
admission policy; prime\_math's first path contains the same power filters.
In a bounded check of the pinned extraction and normalisation helpers, the
candidate difference strings for all four heavy pairs were rejected by both
paths' admission gates. The check ran no symbolic evaluation or complete
grading path; early string decisions and other inputs are outside it. The
filters also reject ordinary expressions such as $2^{20}$, so this is a
compute-versus-admission trade-off, not evidence of superior accuracy.

\begin{table}[t]
\centering\small
\begin{tabular}{p{0.12\linewidth}p{0.41\linewidth}p{0.37\linewidth}}
\toprule
Path & Time and execution boundary & Failure information \\
\midrule
TRL & Calling thread; 10\,s parse / 5\,s compare on main thread, disabled otherwise & Unparseable reference: None; comparison timeout: $0$ \\
Ray example & Calling process; both library limits set to zero & Same reward conversion; no library deadline \\
verl main & Math-Verify is opt-in; prime\_math has subprocess decoration plus the incompatible fallback above & Caught exceptions become false; no failure flag on those paths \\
OpenRLHF & Power filters; 120\,s alarm on main thread, disabled otherwise & Filter refusal, timeout and other exceptions become false \\
\bottomrule
\end{tabular}
\caption{Pinned code comparison for the specified paths. The bounded
fallback and admission-gate checks were executed; full framework configurations
were not.}
\label{tab:frameworks}
\end{table}

\section{Conditional changes in GRPO advantages}
\label{sec:update}
This section is an algebraic illustration, not a measured learning effect
or a new theory of verifier noise. For a group of $G=8$ binary rewards,
fix the sample standard deviation and epsilon as
\[
 \bar r=G^{-1}\sum_i r_i,\qquad
 s^2=(G-1)^{-1}\sum_i(r_i-\bar r)^2,\qquad
 A_i=(r_i-\bar r)/(s+10^{-6}).
\]
Assume the first answer would receive one if it were decided, but its
observed undecided verdict is replaced by zero. Table~\ref{tab:advantages}
shows the advantage of that answer. Constant groups have $s=0$ and all
advantages zero under this convention.

\begin{table}[h]
\centering\small
\begin{tabular}{@{}lrr@{}}
\toprule
Complete clean reward vector & Before zero replacement & After zero replacement \\
\midrule
$(1,1,0,0,0,0,0,0)$ & $1.620182$ & $-0.353552$ \\
$(1,0,0,0,0,0,0,0)$ & $2.474867$ & $0$ \\
$(1,1,1,1,1,1,1,1)$ & $0$ & $-2.474867$ \\
\bottomrule
\end{tabular}
\caption{Conditional advantage changes for the first answer, at fixed
standard-deviation and epsilon conventions. No vector is inferred from
the observed marginal failure rate.}
\label{tab:advantages}
\end{table}

A positive advantage can become negative or zero; a previously zero
advantage can become negative. The group composition matters. In the first
example, the remaining positive answer's advantage changes from 1.620182
to 2.474867. If instead the undecided row is excluded from \emph{both}
group moments and the loss, the remaining positive answer has advantage
2.267781, and the six negative answers each have $-0.377963$.
This masking policy also changes the loss denominator: with equal response
lengths and a token-mean loss over active answers, it is proportional to
seven rather than eight responses. It is not equivalent to assigning the
failed answer advantage zero after computing moments over all eight.
The all-negative and all-positive remaining groups are constant and have
zero advantages. These policies require an observable failure indicator;
they do not supply the unknown clean reward, restore its learning signal,
or establish a better trained model. Existing work already treats noisy
rewards, degenerate groups and normalisation
\cite{cai2025noisy,nie2026starvation,yang2026biased}.
Plesner et al.\ \cite{plesner2026imperfect} also distinguish advantage
changes from clipped gradients and optimisation trajectories. The examples
here illustrate conditional reward handling, not an optimizer update.

%% ---------------------------------------------------------------------------
\section{Implications for reward-pipeline design}
\label{sec:implications}

\paragraph{Verify the contract at the consumer boundary.}
The wrapper, manager and aggregator expose different views of the same
attempt. A manager-side flag is useful only if the intended consumer can
recover its row identity and meaning. The observed position-dependent
outcomes justify testing sparse \texttt{error} and \texttt{timeout} keys at
every row position, across mixed-success chunks and padding/trimming
boundaries. Our local prototype supplies an aggregation regression target;
the consumer still needs to define whether \texttt{None} means absent,
unknown or a value excluded from statistics.

\paragraph{Treat a timeout change as a change in execution policy.}
The main replay contains rare limit-sensitive verdicts, but the shared-pool
experiments show a separate consequence of extending a deadline. When an
outer wait expires first, returning a score need not stop the child work;
killing a shared pool can affect concurrent requests. Default-limit
controls, serial calls and the longer-outer-wait controls distinguish these
effects in the tested configuration. A deadline adjustment should therefore
be checked for both verdict changes and effects on neighbouring requests,
with resource interventions recorded separately.

\paragraph{Separate observability from a claim about learning.}
Retaining a failure flag establishes what happened to a scoring attempt.
It does not establish the mathematical answer, determine the correct
exclusion policy or demonstrate improved training. The labelled replay
also shows that answer representation can alter disagreement much more
than increasing time limits in that selected set. A subsequent training
evaluation would need to fix extraction, adjudicate disputed answers, and
compare explicitly defined inclusion and normalisation policies. The
current evidence provides concrete cases and regression requirements for
that evaluation.

\section{Related work}
\label{sec:related}

\paragraph{Verifier errors in RLVR.}
Cai et al.\ \cite{cai2025noisy} model an imperfect verifier as a noisy reward
channel with false-negative and false-positive rates and give population-level
results and corrections; a fixed broken scoring worker under uniform random
assignment reduces to their channel, but their results are stated before
GRPO's finite-sample normalisation, so they do not transfer to the default
update unchanged. Egashira et al.\ \cite{egashira2026delay} inject systematic verification
errors and observe delay, plateau or collapse; Plesner et al.\
\cite{plesner2026imperfect} report modest mean performance losses in their
multi-seed code experiments with resampled reward noise, while their
uncertainty intervals allow larger losses; Zhang \cite{zhang2026leaky} contrasts training on
code test suites with and without natural false positives. These works
model or measure the effects of erroneous rewards. Our measurements locate
the operations and aggregation policies that turn an unavailable decision
into a reward or a batch exception.

\paragraph{Auditing verifiers and harnesses.}
Xin \cite{xin2026wherefails} audits four verifiers with metamorphic tests and
classifies verdicts as true, false, timeout, error and crash; a follow-up
\cite{xin2026groupcorr} finds that verifier errors are correlated within GRPO
groups and releases the human-labelled rollouts we reuse. Its discussion
separates pooled dependence from prompt difficulty; our finite-corpus
co-occurrence description retains that distinction. Verifier fuzzing \cite{ray2026fuzzing}, allocation and routing of repeated
or multiple verifier calls \cite{hu2026vstress,hu2026routeaudit},
reliability-first advantages \cite{tian2026stargrpo} and harness checks
\cite{pan2026probe} address neighbouring questions. Rollout Cards
\cite{masters2026rolloutcards} includes a source audit of failure recording
in 50 repositories and documents exception-to-zero coercion at aggregation
boundaries. This motivates checking the consumer boundary as well as the
scorer. Our executed trace isolates a row-position effect in a particular
manager and both of its aggregation paths; the 21 fixtures turn that trace
into a concrete regression requirement.

\paragraph{The closest comparisons and the added evidence.}
The closest code-level records are \#6545 and \#7151
\cite{verl6545,verl7151}. The latter also considers multiple workers and
the fully asynchronous return path. Our study does not claim priority for
the first-row-key mechanism, union-and-fill repair or generic cross-worker
alignment. It adds a trace from the specific \#8015 worker failure through
the manager's sparse flag to returned rewards, together with fixed-response
measurements and matched timeout/recovery controls. The local prototype
serves as a regression specification, not as a more complete framework fix.

Xin's verifier audit~\cite{xin2026wherefails} uses certified answer
transformations to distinguish returned verdicts from execution failures.
We reuse the distinction while following runtime events through a wrapper,
manager and consumer. Rollout Cards~\cite{masters2026rolloutcards} already
establishes that failure recording can disappear across layers; the
row-position mechanism identifies the specific operation responsible in
the tested configuration. Finally, the concurrent reset interaction was
reported in the upstream discussion~\cite{verl8015discussion}. The natural
heavy-input replay adds its manager-level consequence and matched default,
serial, single-heavy and longer-wait controls. These comparisons delimit
the added evidence without treating a known failure class as a new concept.

\paragraph{Failures in RL infrastructure.}
Zhang et al.\ \cite{zhang2026rftfailure} inject training faults and detect
them from training curves. iCoder-27B \cite{yang2026icoder} separates
infrastructure from policy errors and normalises over valid samples. On the
optimisation side, Nie et al.\ \cite{nie2026starvation} and Yang et al.\
\cite{yang2026biased} analyse degenerate groups and the bias of group-relative
advantages; any statement about how undecided zeros shift degenerate-group
rates must account for both. The pinned advantage examples here show
conditional changes in finite reward vectors; they neither estimate the
frequency of those vectors in training nor reproduce the theoretical
results of these studies.

%% ---------------------------------------------------------------------------
\section{Limitations}
\label{sec:limitations}

\paragraph{Execution scope.}
The measurements use one Apple-silicon macOS machine. Timing, signal
delivery and process start methods can differ on other systems.
Generation and Ray transport were substituted; real scoring, manager
and aggregation functions were exercised, but no complete trainer was
run. The Math-Verify wrapper is optional in verl's default routing, and
its adoption rate is unknown. The metadata prototype covers local batch
assembly; downstream handling of missing values needs a separate check.

\paragraph{Data and reference validity.}
The main corpus contains responses already filtered by a verifier or
judge, with one answer sampled per problem. Its replay rates do not
estimate training-time incidence. The longer-limit reference is the
same verifier, and the separate human-labelled set is stratified,
contains extracted answers, and was not independently re-adjudicated.
Neither longer limits nor answer formatting establishes mathematical
truth. Finite-sample recovery bounds and binomial working-model bounds
answer different questions; the latter add homogeneity and independence
assumptions not guaranteed by the corpus.

\paragraph{Supplementary analysis.}
The all-answer extension combines the audited main counts with later
complement replays under a newer resource policy. Within-question
co-occurrence uses exact problem text, not semantic deduplication, and
groups have one to six answers rather than the eight-answer training
groups used in the conditional calculations. Natural heavy-input scenes
are selected diagnostics with explicitly changed limits. Their matched
controls identify behaviour in those scenes rather than its frequency
in unconstrained execution.

\paragraph{Learning effects.}
The advantage examples use declared reward vectors, variance conventions
and a numerical stabiliser. No policy gradients, training trajectories,
or model-performance changes were measured. Observability is a software
contract property; whether an exclusion or correction policy improves
learning remains a separate empirical question.

\section{Conclusion}

At the tested verl head, the wrapper's failure-reporting policy composes
poorly with the already reported first-row schema assumption: preserving
an error at the manager boundary does not ensure that it survives batch
assembly. Its row position can determine whether the flag disappears or
assembly fails. The 21 fixtures make reward preservation and metadata
alignment explicit regression requirements; deployment still requires a
defined missing-value policy.

The fixed response corpora show rare limit-sensitive changes and leave
some heavy inputs undecided. Matched replay separates a returning timeout
from unfinished background computation, and connects natural heavy inputs
to an already reported reset interaction under relaxed diagnostic limits.
These results support testing failures across the wrapper, manager and
aggregation boundaries together. They do not establish training incidence
or a learning benefit from a particular handling policy.

\appendix
\section{Reproduction}

The companion archive, \path{rewardfail-companion.zip}, contains the source,
compact reports, derived score records and selected event traces.
Its \path{package_manifest.json} verifies the distributed files; the
original-fingerprint index described below refers to the separately
retained complete records. Public copies substitute portable path labels
without changing reported scores or counts. The archive's README explains
which checks run without the omitted datasets and full raw logs.
The companion source pins the full dependency list in
\path{requirements.lock}, including Math-Verify 0.9.0 and SymPy 1.14.0.
The tested wrapper and manager are at verl pull-request head
\texttt{0c4c3d4}; the comparison wrapper and prime\_math audit are at
\texttt{acad987}. OpenR1-Math-220k is at revision
\path{e4e141ec9dea9f8326f4d347be56105859b2bd68}, default configuration.
The human-label CSV revision and provenance are cited in
\cite{humanlabels}. Its pinned repository carries an MIT licence, and the
pinned OpenR1 data card declares Apache-2.0. These are upstream licence
statements, not a separate audit of all source-data rights. The following
records are kept separately:
\begin{itemize}[leftmargin=*]
\item Injection mechanisms: \path{results/c1c/run1/} and
      \path{results/c1c/run2/}; original execution \texttt{cc87848},
      with two recorded runs compared by \path{experiments/compare_runs.py}.
\item Formal main measurement: \path{results/s2_main/report.json} and
      raw directory \path{main-20261007-202545}. Subsequent experiments
      never change this directory, the formal report or its thresholds.
\item Supplementary diagnostics: \path{results/s2_u18/report.json},
      raw directory \path{u18-20261008-104342}.
\item Labelled-answer cross-check: \path{results/s2_e3/report.json},
      raw directory \path{e3-20261008-codex}; driver at \texttt{2df9212}.
\item All-answer extension: \path{results/s2_e2/report.json}, raw directory
      \path{e2-complement-20261008-codex}; driver at \texttt{4aa71b9}.
      The 100,034 complement answers use the original Parquet reader,
      separate answer-index identifiers and the original worker unchanged.
\item Shared-pool controls and scenes: \path{results/s2_e1/selftest.json},
      \path{results/s2_e1/report.json}, \path{report_eight.json} and
      \path{report_repeat.json} within that directory. Raw directories are
      \path{e1-core-20261008-codex}, \path{e1-eight-20261008-codex} and
      \path{e1-repeat-20261008-codex}; scene snapshots are \texttt{8b8321e},
      \texttt{c026164} and \texttt{3bd7f94}, respectively. The committed
      reports are compact indexes with hashes of the complete raw records;
      they do not duplicate all fixed input strings.
\item Sequential-parent and longer-outer controls:
      \path{results/s2_e1/report_controls.json}, raw directory
      \path{e1-controls-20261008-codex}; scene enumeration at
      \texttt{ca0503b}, execution identical to the earlier scene runner.
\item Local batch-assembly proposal: \path{results/s2_e1/aggregation_contract.json},
      generated by \path{experiments/s2_contract.py} at \texttt{ca0503b}.
      Pinned string-preprocessing checks:
      \path{results/s2_e1/normalization_filter.json}; these do not execute
      SymPy or a complete alternate verifier.
\end{itemize}

Each run records per-file SHA-256 code/data locks, specifications, child
and parent events, pressure readings and intervention/restart history.
Supplementary drivers use the locked worker's existing inline-text path;
they do not edit the main worker or upstream wrapper files. Replay commands
use fresh run directories and are launched with \texttt{nohup}, for example
\begin{verbatim}
nohup .venv/bin/python -I -B experiments/s2_e3.py \
  --run-dir <STAGE2>/runs/e3-new > <STAGE2>/runs/e3-new.log 2>&1 &
nohup .venv/bin/python -I -B experiments/s2_e1.py \
  --run-dir <STAGE2>/runs/e1-new > <STAGE2>/runs/e1-new.log 2>&1 &
nohup .venv/bin/python -I -B experiments/s2_e2.py \
  --run-dir <STAGE2>/runs/e2-new > <STAGE2>/runs/e2-new.log 2>&1 &
\end{verbatim}
Run one measurement controller at a time. Green or warning memory pressure
does not impose a static cap; critical pressure pauses/stops the run and
is recorded as a local intervention, not as a natural verifier failure.
The new drivers' own synthetic checks and raw-record reconstructions were
recorded with their code/source locks. Reassembly of the new raw parent and
child events and a separate count reconstruction matched the reported
results; these checks use the same machine and underlying verifier.

The relative-path index \path{results/artifact_manifest.json} records
SHA-256 hashes and byte sizes for 1,572 files, eight complete raw-run
directories and 32 frozen report values. Its six configurable roots locate
the companion repository, raw stage-two records, separate audits, data and
the two upstream checkouts. The read-only
\path{experiments/check_artifacts.py} verifies this inventory without
scoring answers. Complete, missing-root and corrupt-index checks returned
the declared pass, partial and fail statuses. File relocation was tested
in temporary-directory fixtures; mapping the full inventory to different
root aliases was tested on this host. Neither check is an independent
experimental replication on another machine. Historical drivers and
memory-pressure monitoring still depend on the original macOS layout.

The original main run's raw report is its historical execution report;
\path{results/s2_main/report.json} is the audited reanalysis. Executed
code snapshots and their locks take precedence over current source when
interpreting a run. The inventory does not select redistribution licences
or confer rights to upstream data or software. Source status was refreshed
on 9 October 2026; the added sparse-metadata pull-request records were
checked on 12 October 2026. Cited-paper metadata was checked on
8 October 2026.
Original companion code is licensed under Apache-2.0; this does not replace
upstream licences. This manuscript is licensed under the Creative Commons
Attribution 4.0 International licence
(\href{https://creativecommons.org/licenses/by/4.0/}{CC BY 4.0}). Substantial AI assistance is
disclosed in the author contribution statement.

\paragraph{Conflicts of Interest.} The author declares no conflicts of interest.

\bibliographystyle{plainnat}
\begingroup
\raggedright
\Urlmuskip=0mu
\bibliography{refs}
\endgroup

\end{document}
